\documentclass[11pt,a4paper]{article}
\usepackage[utf8]{inputenc}
\usepackage[T1]{fontenc}
\usepackage{amsmath,amssymb,amsthm}
\usepackage{mathtools}
\usepackage{listings}
\usepackage{xcolor}
\usepackage{booktabs}
\usepackage{hyperref}
\usepackage{geometry}
\usepackage{enumitem}
\usepackage{graphicx}
\usepackage{url}

\geometry{margin=1in}

% Theorem environments
\theoremstyle{definition}
\newtheorem{definition}{Definition}[section]
\newtheorem{theorem}{Theorem}[section]
\newtheorem{corollary}{Corollary}[theorem]
\newtheorem{lemma}{Lemma}[section]
\newtheorem{conjecture}{Conjecture}[section]

% Math operators
\DeclareMathOperator{\I}{I}
\newcommand{\Entropy}{\mathrm{H}}
\newcommand{\given}{\mid}
\newcommand{\indep}{\perp}

% Code listing style
\lstset{
    basicstyle=\ttfamily\small,
    keywordstyle=\color{blue},
    commentstyle=\color{gray},
    stringstyle=\color{red},
    numbers=left,
    numberstyle=\tiny\color{gray},
    stepnumber=1,
    numbersep=5pt,
    frame=single,
    breaklines=true,
    showstringspaces=false
}

% Title information
\title{Dual Privacy Architecture: Information-Theoretic Bounds on Agent Reconstruction\\
\large Mathematical Framework for Swordsman-Mage Separation}
\author{privacymage}
\date{November 25, 2025\\Version 3.2}

\begin{document}

\maketitle

\begin{abstract}
We introduce the Swordsman and Mage as fundamental privacy primitives for dual-agent architectures, establishing rigorous information-theoretic bounds when conditional independence $(Y_S \indep Y_M) \given X$ is enforced between these agents' observations. The Swordsman (S) controls privacy boundaries through selective measurement, while the Mage (M) projects delegated agency using only S-authorized observations.

\textbf{Proven Results:} We prove that this separation enables an additive bound on mutual information: $\I(X; Y_S, Y_M) \leq \I(X; Y_S) + \I(X; Y_M)$. Combined with budget constraints $C_S + C_M < \Entropy(X)$, this establishes a reconstruction ceiling $R_{\max} < 1$ that no adversary can exceed regardless of computational resources. Via Fano's inequality, we establish a fundamental error floor: $P_e \geq 1 - (\I(X; Y) + 1)/\Entropy(X)$, guaranteeing minimum reconstruction error when $R_{\max} < 1$. We further prove graceful degradation under approximate separation.

\textbf{Implementation Framework:} We provide practical budget estimation methods, isolation verification protocols, and side-channel resistance models based on covert channel analysis. We integrate zero-knowledge proof systems for cryptographic enforcement of separation and budget compliance, providing concrete constructions using Groth16, PLONK, and Nova protocols.

\textbf{Theoretical Predictions:} We present theoretical conjectures about potential optimal allocation patterns, including a golden ratio hypothesis ($\phi \approx 1.618$) and tetrahedral emergence properties. These remain unproven mathematical conjectures requiring both formal proof and empirical validation. We outline specific research directions and invite collaboration on implementation and validation.
\end{abstract}

\section*{Nature of This Work}

\textbf{What Is Proven:} The core information-theoretic results (additive bounds under separation, reconstruction ceilings, error floors) are rigorously proven using established information theory.

\textbf{What Is Theoretical:} The golden ratio optimization hypothesis and tetrahedral emergence predictions are unproven mathematical conjectures. They represent interesting theoretical possibilities but have not been formally derived from first principles.

\textbf{What Is Missing:} No implementations exist. No empirical data has been collected. No observations have been made. This is purely theoretical and mathematical work at present.

\textbf{What We Seek:} Collaboration from theorists to prove or disprove the conjectures, and from practitioners to build implementations and collect empirical data.

\section{Introduction}

\subsection{Motivation}

The deployment of autonomous AI agents acting on behalf of humans creates a fundamental tension: agents require information about their principals to act effectively (delegation), yet this same information enables reconstruction of sensitive behavioral patterns (privacy loss). Traditional single-agent architectures cannot resolve this tension---the same system handling both functions creates an inherent conflict of interest.

We propose the Swordsman and Mage as dual privacy primitives that resolve this tension through architectural separation:

\begin{itemize}
    \item \textbf{The Swordsman (S):} A privacy-enforcement primitive that controls information disclosure through selective measurement, acting as the guardian of boundaries
    \item \textbf{The Mage (M):} A delegation primitive that projects agency into the world using only Swordsman-authorized observations, acting as the executor of capabilities
\end{itemize}

These primitives are not mere metaphors but formal architectural components with precise information-theoretic properties. The key insight: enforcing conditional independence between the Swordsman and Mage observations creates provable reconstruction bounds.

\subsection{Contributions}

\textbf{Proven Results:}

\begin{itemize}
    \item \textbf{Separation Lemma (Theorem 3.1):} Under $(Y_S \indep Y_M) \given X$, mutual information becomes additive
    \item \textbf{Reconstruction Ceiling (Corollary 3.2):} With $C_S + C_M < \Entropy(X)$, reconstruction efficiency $R_{\max} < 1$
    \item \textbf{Error Floor (Theorem 3.3):} Fano's inequality establishes minimum error $P_e \geq 1 - (\I(X;Y) + 1)/\Entropy(X)$
    \item \textbf{Robustness Analysis (Theorem 3.4):} $\varepsilon$-approximate separation degrades bounds gracefully
\end{itemize}

\textbf{Implementation Framework:}

\begin{itemize}
    \item Practical budget estimation and monitoring methods
    \item Isolation verification and enforcement protocols
    \item Side-channel resistance model with concrete parameters
    \item Zero-knowledge proof constructions for cryptographic separation enforcement
    \item ZKP-based budget compliance verification protocols
    \item Concrete specifications for Groth16, PLONK, and Nova-based implementations
\end{itemize}

\textbf{Theoretical Predictions} (unproven conjectures):

\begin{itemize}
    \item Golden ratio convergence hypothesis in optimal allocations
    \item Tetrahedral emergence hypothesis for system architecture
\end{itemize}

\subsection{Related Work}

This work differs from existing privacy frameworks:

\begin{itemize}
    \item \textbf{Differential Privacy} \cite{dwork2014}: Adds calibrated noise; we enforce structural separation
    \item \textbf{Secure Multi-Party Computation} \cite{goldreich2004}: Distributes computation; we distribute observation rights
    \item \textbf{Information Flow Control} \cite{sabelfeld2003}: Tracks taint; we bound reconstruction
    \item \textbf{Covert Channel Analysis} \cite{millen1987}: Models side-channels; we apply to privacy architectures
    \item \textbf{Zero-Knowledge Proofs} \cite{groth2016,gabizon2019}: Verifiable computation; we apply to privacy budget enforcement
\end{itemize}

\section{Model and Preliminaries}

\subsection{Basic Framework}

Let $X$ be a secret over finite alphabet $\mathcal{X}$ with $\Entropy(X) > 0$. Two agents produce observations:
\begin{align}
Y_S &= E_S(X, N_S) \\
Y_M &= E_M(X, N_M)
\end{align}
where $N_S, N_M$ are independent local randomness sources.

\subsection{The Swordsman and Mage Primitives}

\begin{definition}[Swordsman Primitive]
The Swordsman $S$ is a privacy-enforcement agent characterized by:
\begin{itemize}
    \item Measurement function $E_S$ that implements selective disclosure
    \item Information budget $C_S$ controlling maximum leakage: $\I(X; Y_S) \leq C_S$
    \item Primary objective: minimize reconstruction while enabling necessary delegation
\end{itemize}
\end{definition}

\begin{definition}[Mage Primitive]
The Mage $M$ is a delegation agent characterized by:
\begin{itemize}
    \item Projection function $E_M$ operating on S-authorized information
    \item Information budget $C_M$ for capability execution: $\I(X; Y_M) \leq C_M$
    \item Primary objective: maximize utility under privacy constraints
\end{itemize}
\end{definition}

The critical architectural requirement: $(Y_S \indep Y_M) \given X$ (conditional independence).

\subsection{Formal Definitions}

\begin{definition}[Separation Condition]
The architecture enforces $(Y_S \indep Y_M) \given X$.
\end{definition}

\begin{definition}[Information Budgets]
$\I(X; Y_S) \leq C_S$, $\I(X; Y_M) \leq C_M$.
\end{definition}

\begin{definition}[Reconstruction Efficiency]
$R \triangleq \I(X; Y)/\Entropy(X) \in [0, 1]$.
\end{definition}

\subsection{Threat Model}

\textbf{Assumptions:}

\begin{itemize}
    \item Passive adversary observing $(Y_S, Y_M)$
    \item Separation enforced through architectural boundaries
    \item Known distributions $P(X)$, encoding functions $E_S, E_M$
\end{itemize}

\textbf{Explicitly Out of Scope:}

\begin{itemize}
    \item Active attacks modifying agent behavior
    \item Side-channels on separation mechanism itself
    \item Temporal correlation across sessions
\end{itemize}

\section{Proven Information-Theoretic Results}

\subsection{The Separation Lemma}

\begin{theorem}[Additive Bound Under Separation]\label{thm:separation}
If $(Y_S \indep Y_M) \given X$ holds, then:
\[
\I(X; Y_S, Y_M) \leq \I(X; Y_S) + \I(X; Y_M)
\]
\end{theorem}

\textit{Note:} Equality holds if and only if $Y_S$ and $Y_M$ are additionally marginally independent. Conditional independence $(Y_S \indep Y_M) \given X$ alone permits dependence through the shared cause $X$, yielding the inequality. The inequality suffices for all downstream guarantees.

\textit{Proof:}

By the chain rule for mutual information:
\[
\I(X; Y_S, Y_M) = \I(X; Y_S) + \I(X; Y_M \given Y_S)
\]

Under conditional independence $(Y_S \indep Y_M) \given X$, we have $\I(Y_M; Y_S \given X) = 0$, which implies:
\[
\Entropy(Y_M \given X, Y_S) = \Entropy(Y_M \given X)
\]

Therefore:
\begin{align}
\I(X; Y_M \given Y_S) &= \Entropy(Y_M \given Y_S) - \Entropy(Y_M \given X, Y_S) \\
&= \Entropy(Y_M \given Y_S) - \Entropy(Y_M \given X) \\
&\leq \Entropy(Y_M) - \Entropy(Y_M \given X) \\
&= \I(X; Y_M)
\end{align}

This completes the proof. $\square$

\begin{corollary}[Reconstruction Ceiling]\label{cor:ceiling}
If $C_S + C_M < \Entropy(X)$, then $R_{\max} = (C_S + C_M)/\Entropy(X) < 1$.
\end{corollary}

\textbf{Critical Clarification:} Separation alone is insufficient. Consider $X$ binary with $Y_S = X$ and $Y_M$ independent noise. Then $(Y_S \indep Y_M) \given X$ holds but $R = 1$. The ceiling requires \textbf{BOTH} separation (for additivity) \textbf{AND} budget constraints (for the bound).

\subsection{Error Lower Bound}

\begin{theorem}[Fano-Based Error Floor]\label{thm:fano}
For any estimator $\hat{X}(Y)$ and finite alphabet $\mathcal{X}$:
\[
P_e \triangleq \Pr[\hat{X}(Y) \neq X] \geq \frac{\Entropy(X \given Y) - 1}{\log(|\mathcal{X}| - 1)}
\]

Which can be rewritten using $\I(X; Y) = \Entropy(X) - \Entropy(X \given Y)$ as:
\[
P_e \geq \frac{\Entropy(X) - \I(X; Y) - 1}{\log(|\mathcal{X}| - 1)}
\]

For large alphabets where $\log(|\mathcal{X}| - 1) \approx \Entropy(X)$:
\[
P_e \gtrsim 1 - \frac{\I(X; Y) + 1}{\Entropy(X)} = 1 - R - \frac{1}{\Entropy(X)}
\]
\end{theorem}

\textit{Proof:}

Apply Fano's inequality in its standard form. For any estimator $\hat{X}(Y)$, the conditional entropy $\Entropy(X \given Y)$ bounds the error probability:
\[
\Entropy(X \given Y) \leq h(P_e) + P_e \cdot \log(|\mathcal{X}| - 1)
\]
where $h(\cdot)$ is binary entropy. Since $h(P_e) \leq 1$, we have:
\[
\Entropy(X \given Y) \leq 1 + P_e \cdot \log(|\mathcal{X}| - 1)
\]

Rearranging:
\[
P_e \geq \frac{\Entropy(X \given Y) - 1}{\log(|\mathcal{X}| - 1)}
\]

Using $\I(X; Y) = \Entropy(X) - \Entropy(X \given Y)$ and $R = \I(X; Y)/\Entropy(X)$:
\begin{align}
P_e &\geq \frac{\Entropy(X) - \I(X; Y) - 1}{\log(|\mathcal{X}| - 1)} \\
&\geq 1 - \frac{\I(X; Y) + 1}{\Entropy(X)} \quad \text{[for large alphabets]} \\
&= 1 - R - \frac{1}{\Entropy(X)}
\end{align}

Therefore, when $R_{\max} < 1$, there exists a fundamental error floor that no estimator can overcome. $\square$

\textbf{Interpretation:} This establishes that when reconstruction efficiency $R_{\max} < 1$ (due to our budget constraints), any adversary must have error probability at least $P_e \geq 1 - R_{\max} - O(1/\Entropy(X))$. For example, if $R_{\max} = 0.7$, then $P_e \geq 0.3 - O(1/\Entropy(X)) \approx 0.3$ for large entropy.

\subsection{Robustness to Approximate Separation}

\begin{theorem}[$\varepsilon$-Approximate Separation]\label{thm:robustness}
If $\I(Y_S; Y_M \given X) \leq \varepsilon$ (approximate separation), then:
\[
\I(X; Y_S, Y_M) \leq \I(X; Y_S) + \I(X; Y_M) + \varepsilon
\]
\end{theorem}

\textit{Proof:}

Following the chain rule decomposition:
\[
\I(X; Y_M \given Y_S) = \Entropy(Y_M \given Y_S) - \Entropy(Y_M \given X, Y_S)
\]

With approximate independence:
\[
\Entropy(Y_M \given X, Y_S) \geq \Entropy(Y_M \given X) - \I(Y_S; Y_M \given X) \geq \Entropy(Y_M \given X) - \varepsilon
\]

Therefore:
\[
\I(X; Y_M \given Y_S) \leq \Entropy(Y_M) - \Entropy(Y_M \given X) + \varepsilon = \I(X; Y_M) + \varepsilon
\]

Combining with $\I(X; Y_S, Y_M) = \I(X; Y_S) + \I(X; Y_M \given Y_S)$ completes the proof. $\square$

This shows graceful degradation: small violations of separation cause proportionally small increases in leakage.

\section{Side-Channel Analysis and Robustness}

\subsection{Connection to Covert Channel Theory}

Our analysis builds on established covert channel capacity results. For $d$ side-channel observations, we model reconstruction growth as:
\[
R(d) = R_{\max} \cdot \frac{\ln(1 + d/d_0)}{\ln(1 + d_{\max}/d_0)}
\]

\textit{Model Assumption:} This logarithmic functional form is adopted as a modeling assumption based on theoretical parallels, not derived from first principles. Empirical validation is required for specific deployment contexts.

This mirrors:

\begin{itemize}
    \item \textbf{Shannon capacity} of noisy channels \cite{cover2006}
    \item \textbf{Timing channel capacity} under sampling \cite{giles2002}
    \item \textbf{Power analysis leakage} with noise \cite{kocher1999}
\end{itemize}

\subsection{Justification for Logarithmic Growth}

The logarithmic form emerges from:

\begin{itemize}
    \item \textbf{Information-theoretic argument:} Diminishing information per sample due to:
    \begin{itemize}
        \item Temporal correlation in behavioral patterns
        \item Finite substrate entropy $\Entropy(X)$
        \item Measurement quantization effects
    \end{itemize}
    \item \textbf{Theoretical foundation:} Previous theoretical work shows logarithmic scaling in:
    \begin{itemize}
        \item Cache timing attacks \cite{yarom2014}
        \item Power analysis \cite{mangard2007}
        \item Traffic correlation \cite{murdoch2005}
    \end{itemize}
\end{itemize}

\subsection{Validation Methodology}

To verify theoretical guarantees in practice:

\begin{itemize}
    \item \textbf{Simulation Framework:}
    \begin{itemize}
        \item Generate test distributions with known $\Entropy(X)$
        \item Implement $S$ and $M$ with controlled coupling
        \item Measure actual $\I(X; Y_S, Y_M)$ vs theoretical bounds
    \end{itemize}
    \item \textbf{Violation Detection Protocol:}
    \begin{itemize}
        \item Monitor $\I(Y_S; Y_M \given X)$ in real-time
        \item Flag when coupling exceeds threshold $\varepsilon$
        \item Trigger corrective isolation measures
    \end{itemize}
\end{itemize}

\subsection{Parameter Estimation}

For practical systems:

\begin{itemize}
    \item $d_0$: Single-query leakage baseline (empirically measured)
    \item $a_S, a_M$: Channel-specific leakage rates (from profiling)
    \item $d_{\max}$: Threat model bound (e.g., $10^6$ queries/day)
\end{itemize}

\section{Design Principles and Budget Management}

\subsection{The Allocation Problem}

Given total budget $C_T = C_S + C_M < \Entropy(X)$, how should we allocate between privacy protection and delegation capacity?

\subsection{Practical Budget Estimation}

\textbf{Monte Carlo Estimation:}

\begin{lstlisting}[language=Python, caption=Budget estimation]
# Budget estimation
def estimate_mutual_info(samples_X, samples_Y):
    # Use MINE or InfoNCE estimators
    # Return confidence interval
    return I_estimate, confidence_bounds
\end{lstlisting}

\textbf{Runtime Monitoring:}

\begin{itemize}
    \item Track cumulative information release
    \item Implement privacy ledger similar to differential privacy
    \item Trigger alerts when approaching budget limits
\end{itemize}

\subsection{Adaptive Control Framework}

\begin{lstlisting}[language=Python, caption=Adaptive budget controller]
# Adaptive budget controller
class AdaptiveBudgetController:
    def __init__(self, C_S_max, C_M_max):
        self.budget_S = C_S_max
        self.budget_M = C_M_max
        
    def adjust_granularity(self, current_usage):
        if current_usage > 0.8 * self.budget_S:
            # Reduce disclosure precision
            return reduced_precision_mode
\end{lstlisting}

\subsection{Deployment Considerations}

\textbf{Handling Feedback Loops:}

\begin{itemize}
    \item Implement causal isolation barriers
    \item Use temporal batching to break correlations
    \item Apply differential privacy within $S$'s operations when needed
\end{itemize}

\textbf{Active Adversary Mitigations:}

\begin{itemize}
    \item Rate limiting on queries
    \item Anomaly detection for unusual patterns
    \item Cryptographic commitments to prevent manipulation
\end{itemize}

\section{Architectural Implementation}

\subsection{Measurement Design Principles}

\textbf{For Privacy Agent (S):}

\begin{itemize}
    \item Differential privacy within budget $C_S$
    \item Verifiable predicates using zero-knowledge proofs
    \item $k$-anonymity for categorical data
    \item Secure sketches for approximate matching
\end{itemize}

\subsubsection{ZKP-Enhanced Selective Disclosure}

The Swordsman's selective disclosure can be implemented using modern zero-knowledge proof systems:

\begin{itemize}
    \item \textbf{Groth16:} For predicate verification with $O(1)$ proof size
    \begin{itemize}
        \item Constant-size proofs (3 group elements)
        \item Fast verification ($\sim$2ms)
        \item Requires trusted setup per circuit
    \end{itemize}
    \item \textbf{PLONK:} For circuit-based privacy constraints
    \begin{itemize}
        \item Universal trusted setup
        \item Flexible constraint systems
        \item Larger proofs but better composability
    \end{itemize}
    \item \textbf{Nova:} For recursive composition of privacy checks
    \begin{itemize}
        \item No trusted setup
        \item Efficient recursive proof composition
        \item Enables cumulative budget verification
    \end{itemize}
\end{itemize}

\textbf{Concrete Construction:} For attribute disclosure, the Swordsman generates:
\[
\pi_{\text{attr}} = \text{ZKP}\{y_S = E_S(x) \land f(y_S) = 1\}
\]
where $f$ is a public predicate (e.g., ``age $\geq 18$'') and the proof reveals nothing beyond the predicate's truth value.

\textbf{For Delegation Agent (M):}

\begin{itemize}
    \item Minimize observation precision
    \item Batch queries to amortize leakage
    \item Caching to avoid redundant observations
    \item Query rate limiting
\end{itemize}

\subsection{Separation Enforcement}

\textbf{Hardware-Based:}
\begin{verbatim}
+-------------+     Rate-Limited      +-------------+
|   TEE (S)   |<------Channel-------->|   TEE (M)   |
|  Privacy    |                       | Delegation  |
+-------------+                       +-------------+
\end{verbatim}

\textbf{Software-Based:}

\begin{itemize}
    \item Container isolation with syscall filtering
    \item Formal verification of information flows
    \item Capability-based access control
    \item Audit logging for violations
\end{itemize}

\subsubsection{Cryptographic Separation Verification}

Rather than relying solely on trusted isolation, we can prove separation cryptographically:
\[
\pi_{\text{sep}} = \text{ZKP}\{(Y_S \indep Y_M) \given X\}
\]

Using techniques from zero-knowledge proof systems:

\begin{itemize}
    \item Commit to both observations: $c_S = \text{Commit}(Y_S, r_S)$, $c_M = \text{Commit}(Y_M, r_M)$
    \item Prove conditional independence via joint distribution commitments
    \item Enable third-party verification of architectural compliance
\end{itemize}

This approach enables:

\begin{itemize}
    \item \textbf{Auditability:} External parties can verify separation without observing raw data
    \item \textbf{Composability:} Proofs can be recursively composed for multi-agent systems
    \item \textbf{Trustless Enforcement:} No need for trusted hardware or monitors
\end{itemize}

\subsection{Isolation Verification}

\textbf{Formal Verification Approaches:}

\begin{itemize}
    \item Non-interference proofs using tools like FlowDroid
    \item Information flow security verification
    \item Covert channel capacity analysis
\end{itemize}

\textbf{Runtime Instrumentation:}

\begin{itemize}
    \item Monitor timing patterns for correlation
    \item Detect shared resource usage
    \item Track API call sequences
\end{itemize}

\subsubsection{ZKP-Based Budget Compliance}

Zero-knowledge proofs enable verifiable budget tracking without revealing actual observations:

\textbf{Protocol:}

\begin{itemize}
    \item Agent commits to observation: $c_i = \text{Commit}(y_i, r_i)$
    \item Proves cumulative budget compliance:
    \[
    \pi_{\text{budget}} = \text{ZKP}\left\{\sum \I(X; y_i) \leq C_S\right\}
    \]
    \item Verifier checks without learning $\{y_i\}$
\end{itemize}

This addresses the monitoring challenge with cryptographic rather than trusted enforcement.

\subsection{Zero-Knowledge Implementation Patterns}

\subsubsection{Range Proofs for Continuous Variables}

For continuous variables, prove $y_S \in [a,b]$ without revealing exact value using Bulletproofs or similar range proof systems.

\textbf{Example:} Location disclosure
\[
\pi_{\text{loc}} = \text{ZKP}\{\text{distance}(y_S, \text{target}) < 5\text{km}\}
\]

\subsubsection{Set Membership Proofs}

Prove attribute membership in approved sets without revealing which element using Merkle tree commitments and Groth16.

\textbf{Example:} Credential verification
\[
\pi_{\text{cred}} = \text{ZKP}\{y_S \in \text{AccreditedUniversities}\}
\]

\subsubsection{Predicate Verification}

Prove arbitrary predicates $f(Y_S) = 1$ using circuit-based SNARKs (PLONK).

\textbf{Example:} Age verification
\[
\pi_{\text{age}} = \text{ZKP}\{\text{age}(y_S) \geq 18 \land \text{age}(y_S) < 120\}
\]

\subsubsection{Cumulative Budget Tracking}

Maintain cryptographic ledger:
\[
\text{Ledger} = \{(c_i, \pi_i)\} \text{ for } i = 1 \text{ to } t
\]
where $c_i$ commits to observation $i$ and $\pi_i$ proves incremental budget consumption stays below $C_S$.

Using Nova for recursive composition:
\begin{align}
\pi_1 &= \text{ZKP}\{\I(X; y_1) \leq C_S\} \\
\pi_t &= \text{ZKP}\{\pi_{t-1} \land \I(X; y_1, \ldots, y_t) \leq C_S\}
\end{align}

\subsection{Implementation Checklist}

\textbf{Pre-deployment:}

\begin{itemize}
    \item[$\square$] Estimate $\Entropy(X)$ for target domain
    \item[$\square$] Set $C_S + C_M \leq 0.7 \cdot \Entropy(X)$ (safety margin)
    \item[$\square$] Implement separation enforcement
    \item[$\square$] Verify isolation properties
    \item[$\square$] Deploy monitoring infrastructure
\end{itemize}

\textbf{Runtime:}

\begin{itemize}
    \item[$\square$] Track actual $\I(X; Y_S)$ and $\I(X; Y_M)$
    \item[$\square$] Monitor separation violations
    \item[$\square$] Log reconstruction attempts
    \item[$\square$] Adjust budgets adaptively
\end{itemize}

\section{Theoretical Predictions}

\subsection{Golden Ratio Convergence Hypothesis}

\textbf{Theoretical Prediction:} Under certain optimization conditions, dual-agent systems may converge to allocation ratios near the golden ratio $\phi \approx 1.618$.

\textbf{Theoretical Foundation:} Consider value functions of the form $V(C_S, C_M) = U(C_M) \cdot P(C_S) \cdot \Psi(R_{\max})$ where:

\begin{itemize}
    \item $U(C_M)$: Delegation utility (e.g., power-law: $U(x) = x^\alpha$)
    \item $P(C_S)$: Privacy protection value (e.g., $P(x) = x^\beta$)
    \item $\Psi(R_{\max})$: Reconstruction penalty
\end{itemize}

For Lagrangian optimization with power-law utilities, the optimal ratio satisfies:
\[
\frac{\partial V/\partial C_S}{\partial V/\partial C_M} = \frac{\beta}{\alpha}
\]

When $\beta/\alpha \approx \phi$ (protection slightly exceeds delegation), systems may show enhanced stability. This theoretical prediction suggests $\phi$ as a potential optimization attractor under symmetric importance of privacy and utility.

\textbf{Status:} This is a theoretical prediction derived from optimization analysis. Rigorous mathematical proof of universal convergence remains an open problem requiring further validation.

\subsection{Open Research Question: Data Collection Needed}

\textbf{Research Opportunity:} We seek to collect empirical data on $C_S/C_M$ ratio patterns across different implementation domains to test the golden ratio convergence hypothesis. This is purely theoretical at present---no implementations or observations exist yet.

\textbf{Call for Collaboration:} Researchers and developers implementing dual-agent privacy systems are invited to share their allocation patterns, system behaviors, and empirical measurements. This data would be invaluable for validating or refuting the theoretical predictions presented in this paper.

\subsection{Tetrahedral Emergence Hypothesis}

\textbf{Theoretical Prediction:} Sustained $S$-$M$ separation may naturally generate two additional measurement properties, forming a tetrahedral system architecture:

\begin{itemize}
    \item \textbf{Reflect (R):} Temporal accumulation of $S$'s boundary decisions, creating memory properties
    \item \textbf{Connect (C):} Network effects from $M$'s delegation patterns, creating relational properties
\end{itemize}

\textbf{Mathematical Formulation:} If $(Y_S \indep Y_M) \given X$ is maintained over time, the system may generate:
\begin{align}
Y_R &= R(Y_S^{1:t}, \tau) \quad \text{[Memory from S history]} \\
Y_C &= C(Y_M^{1:t}, G) \quad \text{[Network from M interactions]}
\end{align}

By data processing inequality, $\I(X; Y_R) \leq \I(X; Y_S)$ and $\I(X; Y_C) \leq \I(X; Y_M)$, suggesting bounded additional leakage if these properties emerge.

\textbf{Theoretical Rationale:}

\begin{itemize}
    \item \textbf{Minimum for sovereignty:} Two agents create separation, two properties enable persistence
    \item \textbf{Natural dimensions:} Temporal (R) and spatial (C) emerge from operational (S, M)
    \item \textbf{Information conservation:} Additional properties would be bounded by primary agent budgets
\end{itemize}

\textbf{Status:} This is a theoretical prediction about potential system evolution. Whether such emergence is inevitable, optimal, or even desirable requires further mathematical analysis and empirical validation.

\subsection{Testable Predictions}

If tetrahedral emergence occurs in real systems, we would expect to observe:

\begin{itemize}
    \item \textbf{Temporal patterns:} Systems may develop memory/logging behaviors (potential ``Reflect'' property)
    \item \textbf{Network effects:} Inter-agent communication patterns may emerge (potential ``Connect'' property)
\end{itemize}

\textbf{Important:} These are theoretical predictions, not observations. No implementations exist yet to test whether such emergence actually occurs. These predictions provide concrete hypotheses for future experimental work.

\subsection{Future Research Directions}

\begin{itemize}
    \item \textbf{Statistical Validation:}
    \begin{itemize}
        \item Large-scale experiments across domains
        \item Hypothesis testing for allocation patterns
        \item Control for confounding variables
    \end{itemize}
    \item \textbf{Theoretical Investigation:}
    \begin{itemize}
        \item Derive optimal allocation from first principles
        \item Characterize conditions for specific ratio convergence
        \item Prove or disprove emergence hypotheses
    \end{itemize}
    \item \textbf{Practical Extensions:}
    \begin{itemize}
        \item Multi-agent generalizations (beyond two)
        \item Dynamic budget adjustment algorithms
        \item Integration with existing privacy tools
    \end{itemize}
\end{itemize}

\section{Theoretical Predictions and Open Questions}

\textbf{STATUS: PURELY THEORETICAL} --- This section presents unproven mathematical conjectures and theoretical predictions. No implementations, empirical data, or observations exist. These are research hypotheses requiring both formal mathematical proof and experimental validation.

\subsection{Golden Ratio Hypothesis (Unproven)}

\begin{conjecture}[Golden Ratio Optimality - UNPROVEN]\label{conj:golden}
There may exist an optimization principle that drives optimal allocation ratios toward $\phi \approx 1.618$.
\end{conjecture}

\textbf{Theoretical Motivation:}

If a joint optimization objective exists of the form:
\[
\max_{C_S, C_M} U(C_M) \cdot P(C_S, C_M) \quad \text{s.t.} \quad C_S + C_M = B
\]
where $U$ is utility and $P$ is privacy protection, then under certain smoothness and monotonicity conditions, the optimal ratio might be $C_S/C_M = \phi$.

\textbf{What Would Be Required to Validate This:}

\begin{itemize}
    \item Formal mathematical proof deriving $\phi$ from first principles
    \item Precise definitions of utility and privacy functions
    \item Proof of existence and uniqueness of optimal ratio
    \item Characterization of exact conditions under which $\phi$ emerges
    \item Experimental implementations across multiple domains
    \item Statistical validation with large sample sizes
\end{itemize}

\textbf{Current Status:} Pure conjecture. No proof exists. No data exists. This is a mathematical hypothesis awaiting investigation.

\subsection{Budget Allocation: Testable Framework}

\textbf{Proposed Metrics} (for future experimental work):

Define efficiency metrics for implementations:

\begin{itemize}
    \item \textbf{Privacy Efficiency:} $\eta_P = \text{Privacy Achieved} / \text{Budget Spent}$
    \item \textbf{Utility Efficiency:} $\eta_U = \text{Utility Achieved} / \text{Budget Spent}$
    \item \textbf{Joint Efficiency:} $\eta_J = \eta_P \cdot \eta_U$
\end{itemize}

\textbf{Theoretical Predictions} (to be tested):

\begin{itemize}
    \item $\eta_P$ may increase with $C_S$ but with diminishing returns
    \item $\eta_U$ may increase with $C_M$ but with diminishing returns
    \item $\eta_J$ might show local maxima near $C_S/C_M \approx \phi$ if the golden ratio hypothesis holds
\end{itemize}

\textbf{Critical Issues Requiring Resolution:}

\begin{itemize}
    \item ``Privacy Achieved'' and ``Utility Achieved'' need rigorous, measurable definitions
    \item Metrics may be highly implementation-dependent
    \item Optimization landscape likely varies by application domain
    \item No implementations exist yet to test these predictions
\end{itemize}

\section{Tetrahedral Emergence Hypothesis}

\textbf{STATUS: HIGHLY SPECULATIVE} --- This section presents a geometric interpretation that has not been mathematically validated.

\subsection{Four-Force Model}

\begin{conjecture}[Tetrahedral Structure]
The dual-agent system, when operating over time, may exhibit four distinct information flows:

\begin{itemize}
    \item \textbf{Swordsman (S):} Privacy enforcement (boundary-making)
    \item \textbf{Mage (M):} Delegation (capability projection)
    \item \textbf{Reflect (R):} Temporal integration (memory formation)
    \item \textbf{Connect (C):} Network coordination (relationship building)
\end{itemize}
\end{conjecture}

\textbf{Geometric Interpretation:}

If these four forces exist and maintain pairwise relationships, they may form a tetrahedral structure where:

\begin{itemize}
    \item Each vertex represents one force
    \item Each edge represents an information channel
    \item The structure is minimal (no redundant channels)
    \item Stability requires specific geometric relationships
\end{itemize}

\textbf{Mathematical Formulation (Tentative):}

Let $I_{ij}$ denote mutual information between forces $i$ and $j$. The tetrahedron hypothesis suggests:
\begin{align}
I_{SM} &= 0 \quad \text{(proven separation)} \\
I_{SR} &\approx \phi \cdot I_{MR} \quad \text{(conjectured)} \\
I_{SC} &\approx \phi \cdot I_{MC} \quad \text{(conjectured)} \\
I_{RC} &\approx \phi^2 \quad \text{(highly speculative)}
\end{align}

\subsection{Emergence from Dual-Agent Core}

\textbf{Claim 8.1 (UNPROVEN):} Reflect and Connect are not separate agents but emergent properties of sustained $S$-$M$ separation over time.

\textbf{Proposed Mechanism:}

\begin{itemize}
    \item \textbf{Reflect emerges} when Swordsman maintains consistent boundaries across sessions, creating temporal correlation structure
    \item \textbf{Connect emerges} when Mage builds relationships with other agents, creating network correlation structure
\end{itemize}

\textbf{Required Validation:}

\begin{itemize}
    \item Formal definition of ``emergence'' in this context
    \item Proof that $R$ and $C$ are inevitable given $S$-$M$ separation
    \item Demonstration that $R$ and $C$ cannot be reduced to $S$ and $M$
    \item Empirical observation in deployed systems
\end{itemize}

\subsection{Golden Ratio in Tetrahedral Context}

\textbf{Speculation 8.1:} If the four-force tetrahedron exists and maintains stability, the golden ratio may govern relationships:
\[
I_S/I_M = \phi, \quad I_S/I_R = \phi^2, \quad I_S/I_C = \phi^3
\]

\textbf{Geometric Justification (Informal):}

The golden ratio appears in regular geometric structures (pentagons, Fibonacci spirals). If information flows in the tetrahedron seek geometric efficiency, $\phi$-based relationships might provide optimal packing or stability properties.

\textbf{Critical Acknowledgment:} This geometric interpretation is highly speculative and serves primarily as a framework for future investigation. No rigorous derivation exists, and the hypothesis may ultimately prove incorrect.

\section{Discussion}

\subsection{Summary of Proven Guarantees}

We have rigorously established:

\begin{itemize}
    \item Separation enables additive mutual information bounds
    \item Combined with budgets, guarantees $R_{\max} < 1$
    \item Fano's inequality ensures minimum error rates
    \item Approximate separation degrades gracefully
    \item ZKP constructions enable cryptographic enforcement
\end{itemize}

These results hold unconditionally, independent of computational assumptions (except for ZKP implementations which require standard cryptographic hardness).

\subsection{Practical Implementation Pathways}

The framework can be deployed through:

\begin{itemize}
    \item \textbf{Hybrid Approaches:} Combine with differential privacy for additional protection
    \item \textbf{Layered Architecture:} Use TEEs for hardware isolation, containers for software
    \item \textbf{Gradual Adoption:} Start with low-sensitivity applications, expand as confidence grows
\end{itemize}

\subsection{Relationship to Existing Privacy Frameworks}

\begin{table}[h]
\centering
\begin{tabular}{@{}llll@{}}
\toprule
Framework & Focus & Our Approach & Synergy \\
\midrule
Differential Privacy & Statistical noise & Structural separation & Use DP within $S$ \\
Secure MPC & Distributed computation & Distributed observation & Complementary \\
Information Flow Control & Taint tracking & Quantitative bounds & Enhanced metrics \\
\bottomrule
\end{tabular}
\caption{Comparison with existing privacy frameworks}
\end{table}

\subsection{Open Questions and Theoretical Predictions}

\textbf{Unresolved Questions:}

\begin{itemize}
    \item Is two-agent separation optimal for privacy-utility trade-offs?
    \item Can the golden ratio convergence be proven from first principles?
    \item Does tetrahedral emergence occur naturally or require explicit design?
    \item How do these bounds compose in multi-agent systems?
\end{itemize}

\textbf{Theoretical Predictions Requiring Validation:}

\begin{itemize}
    \item Golden ratio proportions as optimization attractors in allocation strategies
    \item Tetrahedral structure as natural consequence of sustained dual-agent separation
    \item Emergent properties $R$ and $C$ as inevitable system evolution
    \item Four-force structure as minimal complete sovereignty architecture
\end{itemize}

\subsection{Limitations and Assumptions}

\textbf{Key Limitations:}

\begin{itemize}
    \item \textbf{Conditional Independence:} Hard to enforce perfectly in practice
    \item \textbf{Passive Adversary:} Active attacks not fully addressed
    \item \textbf{Known Distributions:} Uncertainty in $P(X)$ affects budgets
    \item \textbf{Static Budgets:} Dynamic environments may require adaptation
\end{itemize}

\textbf{Critical Assumptions:}

\begin{itemize}
    \item Architectural isolation is maintainable
    \item Side channels are bounded and detectable
    \item Budget estimation is sufficiently accurate
\end{itemize}

\subsection{Experimental Roadmap}

\textbf{Immediate Next Steps:}

\begin{itemize}
    \item Implement reference architecture with common privacy benchmarks
    \item Develop standardized test suite for separation verification
    \item Create open-source monitoring tools
\end{itemize}

\textbf{Medium-term Goals:}

\begin{itemize}
    \item Deploy in real applications with user studies
    \item Validate across diverse domains
    \item Establish best practices for budget setting
\end{itemize}

\textbf{Long-term Research:}

\begin{itemize}
    \item Prove or refute optimal allocation theorems
    \item Extend to multi-agent and distributed settings
    \item Integrate with emerging privacy technologies
\end{itemize}

\section{Related Extended Work}

\subsection{Connections to Complex Systems}

The potential emergence of structure from simple rules connects to:

\begin{itemize}
    \item \textbf{Self-organization} \cite{prigogine1984}
    \item \textbf{Emergent complexity} \cite{kauffman1993}
    \item \textbf{Scale-free networks} \cite{barabasi2002}
    \item \textbf{Allometric scaling} \cite{west2017}
\end{itemize}

\subsection{Privacy Technology Integration}

Our framework complements:

\begin{itemize}
    \item \textbf{Zero-knowledge proofs:} Can implement $S$'s selective disclosure with cryptographic guarantees (Groth16, PLONK, Nova)
    \item \textbf{Secure enclaves:} Hardware enforcement of separation
    \item \textbf{Homomorphic encryption:} Computation within $M$'s bounds
    \item \textbf{Privacy pools:} Network effects without individual exposure
\end{itemize}

\subsection{Economic and Game Theory}

If golden ratios emerge naturally:

\begin{itemize}
    \item \textbf{Nash equilibria} might converge to $\phi$-proportions
    \item \textbf{Market mechanisms} could discover optimal allocations
    \item \textbf{Evolutionary pressure} might select for tetrahedral structures
\end{itemize}

\textbf{Economic Enforcement of Separation:}

The architectural separation $(Y_S \indep Y_M) \given X$ can be economically enforced through dual-token markets:

\begin{itemize}
    \item SWORD tokens earned exclusively through Swordsman chronicles (privacy domain)
    \item MAGE tokens earned exclusively through Mage chronicles (delegation domain)
    \item Market separation creates economic pressure against agent merger
    \item Budget tracking $C_S$ and $C_M$ verifiable on-chain
    \item Guardian staking (10,000 SWORD) maintains collective compression standards
\end{itemize}

If a First Person attempts to merge agents (violating $(Y_S \indep Y_M) \given X$), they lose earning capability in both domains---creating self-enforcing separation through incentive alignment.

\textbf{Signal-Based Sustainability:}

Rather than speculative token sales, sustainable funding emerges from ceremony and signal fees:

\begin{itemize}
    \item Genesis ceremony: 1 ZEC (\$500 at \$500/ZEC) creates agent pair once per ecosystem
    \item Ongoing signals: 0.01 ZEC (\$5) each, continuous proof-of-comprehension
    \item Fee distribution: 61.8\% transparent pool, 38.2\% shielded pool (internal allocation per ecosystem)
    \item Self-sustaining through activity-based revenue
\end{itemize}

\textbf{Relationship to Proven Results:}

The economic model described above is \textit{one possible implementation}---the mathematical guarantees established in Sections 3--6 hold independent of tokenomic choices. The reconstruction ceiling, error floor, and separation bounds are proven theorems that remain valid regardless of:

\begin{itemize}
    \item Token model (dual, single, or no tokens)
    \item Funding mechanism (ceremonies, grants, or other)
    \item Reward structure (progressive, fixed, or algorithmic)
    \item Network effects (VRCs optional, not required)
\end{itemize}

What the proven mathematics \textit{enables} is sustainable tokenomics grounded in information theory rather than speculation. The budget constraint $C_S + C_M < \Entropy(X)$ provides a scarcity bound; how that bound is implemented economically is a design choice.

\textbf{For complete economic architecture, see companion document:} ``VRC Protocol: Economic Architecture''

\section{Conclusion}

We have established rigorous information-theoretic bounds for dual-agent privacy architectures with enforced separation. The proven results---additive mutual information under separation, reconstruction ceilings below unity, and guaranteed error floors---provide solid foundations for privacy-preserving agent systems.

We integrate zero-knowledge proof systems as core implementation primitives, providing concrete constructions using Groth16, PLONK, and Nova protocols. This enables cryptographic rather than merely architectural enforcement of separation and budget constraints, offering stronger guarantees in adversarial settings.

The key insight remains powerful: structural separation with budget constraints creates fundamental privacy guarantees independent of computational assumptions. This offers a principled approach to the agent privacy problem that complements existing methods.

We present theoretical conjectures about golden ratio optimization and tetrahedral emergence, but emphasize these remain unproven mathematical hypotheses with no empirical validation. The framework's value lies entirely in its proven information-theoretic guarantees and the implementation roadmap it provides for future builders.

\section*{Version Statement}

\textbf{Version 3.2:} This edition clearly separates proven information-theoretic results from unproven theoretical conjectures. The core separation bounds, reconstruction ceilings, and error guarantees are rigorously established through formal proof. Golden ratio hypotheses and tetrahedral emergence remain purely theoretical conjectures with no empirical validation---no implementations or observations exist yet. We explicitly distinguish mathematical theorems from speculative predictions and invite collaboration on both formal proof development and experimental validation. For implementation guidance, see the companion \textit{0xagentprivacy Whitepaper}.

\section*{Acknowledgments}

To the 0xagentprivacy project for exploring privacy-first infrastructure. To BGIN for discussions on AI governance. To Kwaai for advancing personal AI sovereignty. To First Person Project for pioneering user-centric identity frameworks. To the reviewers whose feedback strengthened this work.

\begin{thebibliography}{99}

\bibitem{dwork2014}
Dwork, C. \& Roth, A. (2014). The Algorithmic Foundations of Differential Privacy. \textit{Foundations and Trends in Theoretical Computer Science}.

\bibitem{goldreich2004}
Goldreich, O. (2004). Foundations of Cryptography: Volume 2, Basic Applications. Cambridge University Press.

\bibitem{sabelfeld2003}
Sabelfeld, A. \& Myers, A. C. (2003). Language-based information-flow security. \textit{IEEE Journal on Selected Areas in Communications}, 21(1), 5--19.

\bibitem{kairouz2017}
Kairouz, P., Oh, S., \& Viswanath, P. (2017). The Composition Theorem for Differential Privacy. \textit{IEEE Transactions on Information Theory}, 63(6), 4037--4039.

\bibitem{shannon1948}
Shannon, C. E. (1948). A Mathematical Theory of Communication. \textit{Bell System Technical Journal}.

\bibitem{cover2006}
Cover, T. \& Thomas, J. (2006). Elements of Information Theory. Wiley.

\bibitem{fano1961}
Fano, R. M. (1961). Transmission of Information. MIT Press.

\bibitem{millen1987}
Millen, J. K. (1987). Covert Channel Capacity. \textit{IEEE Symposium on Security and Privacy}.

\bibitem{gray1993}
Gray III, J. W. (1993). On Analyzing the Bus-Contention Channel. \textit{IEEE Computer Security Foundations Workshop}.

\bibitem{giles2002}
Giles, J. \& Hajek, B. (2002). An Information-Theoretic and Game-Theoretic Study of Timing Channels. \textit{IEEE Transactions on Information Theory}, 48(9), 2455--2477.

\bibitem{kocher1999}
Kocher, P., Jaffe, J., \& Jun, B. (1999). Differential Power Analysis. \textit{CRYPTO '99}.

\bibitem{yarom2014}
Yarom, Y. \& Falkner, K. (2014). FLUSH+RELOAD: A High Resolution, Low Noise, L3 Cache Side-Channel Attack. \textit{USENIX Security}.

\bibitem{mangard2007}
Mangard, S., Oswald, E., \& Popp, T. (2007). Power Analysis Attacks: Revealing the Secrets of Smart Cards. Springer.

\bibitem{murdoch2005}
Murdoch, S. J. \& Danezis, G. (2005). Low-Cost Traffic Analysis of Tor. \textit{IEEE Symposium on Security and Privacy}.

\bibitem{dwork2010}
Dwork, C., Rothblum, G. N., \& Vadhan, S. (2010). Boosting and Differential Privacy. \textit{FOCS 2010}.

\bibitem{mironov2017}
Mironov, I. (2017). Rényi Differential Privacy. \textit{IEEE Computer Security Foundations Symposium}.

\bibitem{prigogine1984}
Prigogine, I. \& Stengers, I. (1984). Order Out of Chaos. Bantam Books.

\bibitem{kauffman1993}
Kauffman, S. (1993). The Origins of Order: Self-Organization and Selection in Evolution. Oxford University Press.

\bibitem{barabasi2002}
Barabási, A. L. (2002). Linked: The New Science of Networks. Perseus Books.

\bibitem{west2017}
West, G. (2017). Scale: The Universal Laws of Life and Death in Organisms, Cities and Companies. Penguin.

\bibitem{ball2009}
Ball, P. (2009). Nature's Patterns: A Tapestry in Three Parts. Oxford University Press.

\bibitem{livio2002}
Livio, M. (2002). The Golden Ratio: The Story of Phi. Broadway Books.

\bibitem{groth2016}
Groth, J. (2016). On the Size of Pairing-based Non-interactive Arguments. \textit{EUROCRYPT 2016}. ePrint Archive 2016/260.

\bibitem{gabizon2019}
Gabizon, A., Williamson, Z. J., \& Ciobotaru, O. (2019). PLONK: Permutations over Lagrange-bases for Oecumenical Noninteractive arguments of Knowledge. ePrint Archive 2019/953.

\bibitem{kothapalli2021}
Kothapalli, A., Setty, S., \& Tzialla, I. (2021). Nova: Recursive Zero-Knowledge Arguments from Folding Schemes. ePrint Archive 2021/370.

\end{thebibliography}

\appendix

\section{Proof Details}

\subsection{Complete Chain Rule Expansion}

For completeness, the full chain rule for four variables:
\[
\I(X; Y_S, Y_M, Y_R, Y_C) = \I(X; Y_S) + \I(X; Y_M \given Y_S) + \I(X; Y_R \given Y_S, Y_M) + \I(X; Y_C \given Y_S, Y_M, Y_R)
\]

Each conditional term is bounded by the unconditional under independence assumptions.

\subsection{Fano's Inequality Tightness}

For binary classification with symmetric channels, Fano's bound is nearly tight. For larger alphabets, the bound loosens but remains non-trivial for practical entropy values.

\section{Implementation Pseudocode}

\subsection{Basic Dual-Agent System}

\begin{lstlisting}[language=Python, caption=Basic dual-agent system]
# Basic dual-agent system
class DualAgentPrivacy:
    def __init__(self, entropy_bits, safety_factor=0.7):
        self.H_X = entropy_bits
        self.budget = self.H_X * safety_factor
        
        # Allocation (example values, not empirically validated)
        # Golden ratio hypothesis suggests C_S ~= 0.62, C_M ~= 0.38
        self.C_S = self.budget * 0.62  
        self.C_M = self.budget * 0.38
        
    def enforce_separation(self):
        """Implementation-specific isolation"""
        # Use TEEs, containers, or formal verification
        pass
        
    def measure_leakage(self):
        I_S = self.estimate_mutual_info(self.Y_S, self.X)
        I_M = self.estimate_mutual_info(self.Y_M, self.X)
        I_joint = self.estimate_mutual_info((self.Y_S, self.Y_M), self.X)
        
        # Verify separation
        separation_violation = I_joint - I_S - I_M
        return I_S, I_M, separation_violation
\end{lstlisting}

\subsection{Adaptive Budget Controller}

\begin{lstlisting}[language=Python, caption=Adaptive budget controller]
# Adaptive budget controller
class AdaptiveBudgetController:
    def __init__(self, total_budget, initial_ratio=1.5):
        self.total = total_budget
        self.ratio = initial_ratio
        self.history = []
        
    def update(self, utility_score, privacy_score):
        """Adjust allocation based on performance"""
        self.history.append((utility_score, privacy_score))
        
        if len(self.history) > 100:
            # Analyze trends, adjust ratio
            if utility_score < threshold:
                self.ratio *= 0.9  # Give more to Mage
            elif privacy_score < threshold:
                self.ratio *= 1.1  # Give more to Swordsman
                
    def get_budgets(self):
        C_S = self.total * (self.ratio / (1 + self.ratio))
        C_M = self.total * (1 / (1 + self.ratio))
        return C_S, C_M
\end{lstlisting}

\subsection{Testing Separation Violations}

\begin{lstlisting}[language=Python, caption=Separation violation testing]
# Separation violation testing
def test_separation_violation(system, num_samples=10000):
    """Empirically verify conditional independence"""
    samples = []
    for _ in range(num_samples):
        x = sample_secret()
        y_s, y_m = system.observe(x)
        samples.append((x, y_s, y_m))
    
    # Compute I(Y_S; Y_M | X)
    violation = estimate_conditional_mi(samples)
    
    # Check against threshold
    if violation > epsilon:
        return "VIOLATION DETECTED", violation
    return "SEPARATION MAINTAINED", violation
\end{lstlisting}

\subsection{Budget Compliance Monitoring}

\begin{lstlisting}[language=Python, caption=Budget compliance monitoring]
# Budget compliance monitoring
def monitor_budget_compliance(agent, budget, window=1000):
    """Real-time budget monitoring"""
    cumulative_info = 0
    measurements = deque(maxlen=window)
    
    while True:
        x, y = agent.get_next_observation()
        instant_mi = estimate_instant_mi(x, y)
        measurements.append(instant_mi)
        
        cumulative_info = sum(measurements)
        if cumulative_info > budget:
            agent.trigger_throttling()
            log_budget_violation()
\end{lstlisting}

\end{document}

